% Part:sets-functions-relations
% Chapter: size-of-sets
% Section: pairing-alt

\documentclass[../../../include/open-logic-section]{subfiles}

\begin{document}

\olfileid{sfr}{siz}{pai-alt}

\olsection{একটি বিকল্প যুগলায়ন অপেক্ষক}

\begin{explain}
$\Nat^2$-এর এমন আরও তালিকায়ন আছে, যাদের বিপরীত নির্ণয় করা সহজতর। এখানে তেমন একটি দিচ্ছি। তালিকায়নকে সারণিতে দেখার পরিবর্তে ধনাত্মক পূর্ণসংখ্যার একটি তালিকা দিয়ে শুরু করো, যার সঙ্গে প্রথমে ফাঁকা ঘর যুক্ত আছে। ভাবো, এই ঘরগুলিকে নিচের নিয়মে একের পর এক ক্রমযুগল $\tuple{n,m}$ দিয়ে পূরণ করা হচ্ছে। যেসব ক্রমযুগলের প্রথম স্থানে~$0$ আছে (অর্থাৎ $\tuple{0,m}$), সেগুলি দিয়ে শুরু করো। প্রথমটিকে (অর্থাৎ $\tuple{0,0}$) প্রথম ফাঁকা ঘরে রাখো, তারপর একটি ফাঁকা ঘর ছেড়ে দ্বিতীয়টিকে (অর্থাৎ $\tuple{0,2}$) পরের ফাঁকা ঘরে রাখো, আবার একটি ঘর ছাড়ো, এইভাবে এগোও। আমাদের তালিকায়নের অসম্পূর্ণ শুরুটি এখন এই রকম:
\[\small
\begin{array}{@{}c c c c c c c c c c c@{}}
\mathbf 1 & \mathbf 2 & \mathbf 3 & \mathbf 4 & \mathbf 5 & \mathbf 6 & \mathbf 7 & \mathbf 8 & \mathbf 9 & \mathbf{10} & \dots \\ \\
\tuple{0,0} &  & \tuple{0,1} &  & \tuple{0,2} &  & \tuple{0,3} & & \tuple{0,4} &  & \dots \\
\end{array}
\]
যেসব ঘর এখনও ফাঁকা আছে, সেখানে $\tuple{1,m}$ ক্রমযুগলগুলির জন্য একই কাজ করো; এবারও একটি করে ফাঁকা ঘর ছেড়ে এগোও:
\[\small
\begin{array}{@{}c c c c c c c c c c c@{}}
\mathbf 1 & \mathbf 2 & \mathbf 3 & \mathbf 4 & \mathbf 5 & \mathbf 6 & \mathbf 7 & \mathbf 8 & \mathbf 9 & \mathbf{10} & \dots \\ \\
\tuple{0,0} & \tuple{1,0} & \tuple{0,1} &  & \tuple{0,2} & \tuple{1,1} &
\tuple{0,3} & & \tuple{0,4} &  \tuple{1,2} & \dots \\
\end{array}
\]
একইভাবে $\tuple{2,m}$, $\tuple{3,m}$ ইত্যাদি ক্রমযুগল বসাও। আমাদের সম্পূর্ণ তালিকায়ন তাই এইভাবে শুরু হয়: % OLSIZ-003 TARGET-CORRECTION-BEGIN
\par\smallskip\noindent\textnormal{\small\textbf{উৎস-সংশোধন (OLSIZ-003):} হিমায়িত ইংরেজি উৎসে পরপর সারিগুলি বোঝাতে $\tuple{2,m}$ পরিবারটি দুবার লেখা হয়েছে। অব্যবহিত পরের সারণিতে সারিগুলি $\tuple{2,m}$, $\tuple{3,m}$ ইত্যাদি; বাংলা পাঠে দ্বিতীয় পরিবারটি সেই অনুযায়ী $\tuple{3,m}$ করা হয়েছে।} % OLSIZ-003 TARGET-CORRECTION-END
\[\small
\begin{array}{@{}cc c c c c c c c c c@{}}
\mathbf 1 & \mathbf 2 & \mathbf 3 & \mathbf 4 & \mathbf 5 & \mathbf 6 & \mathbf 7 & \mathbf 8 & \mathbf 9 & \mathbf{10} & \dots \\ \\
\tuple{0,0} & \tuple{1,0} & \tuple{0,1} & \tuple{2,0}  & \tuple{0,2} &
\tuple{1,1} & \tuple{0,3} & \tuple{3,0}  & \tuple{0,4} &  \tuple{1,2} & \dots \\
\end{array}
\]
এই তালিকায়ন অনুযায়ী আগের সারণির ঘরগুলিকে ক্রমসংখ্যা দিলে পরিচ্ছন্ন আঁকাবাঁকা রেখা পাব না; বরং এই বিন্যাস পাব:
\[
\begin{array}{ c | c | c | c | c | c | c | c }
& \mathbf 0 & \mathbf 1 & \mathbf 2 & \mathbf 3 & \mathbf 4 & \mathbf 5 & \dots \\
\hline
\mathbf 0 & 1 & 3 & 5 & 7 & 9 & 11 & \dots \\
\hline
\mathbf 1 & 2 & 6 & 10 & 14 & 18 & \dots & \dots \\
\hline
\mathbf 2 & 4 & 12 & 20 & 28 & \dots & \dots & \dots \\
\hline
\mathbf 3 & 8 & 24 & 40 & \dots & \dots & \dots & \dots \\
\hline
\mathbf 4 & 16 & 48 & \dots & \dots & \dots & \dots & \dots \\
\hline
\mathbf 5 & 32 & \dots & \dots & \dots & \dots & \dots & \dots \\
\hline
\vdots & \vdots & \vdots & \vdots & \vdots & \vdots & \vdots & \ddots\\
\end{array}
\]
দেখা যাচ্ছে,~$0$ নম্বর সারির ক্রমযুগলগুলি তালিকায়নের বিজোড় ক্রমস্থানে আছে, অর্থাৎ ক্রমযুগল $\tuple{0,m}$ আছে $2m+1$ ক্রমস্থানে। দ্বিতীয় সারির ক্রমযুগল $\tuple{1,m}$-গুলি এমন ক্রমস্থানে আছে, যার সংখ্যা একটি বিজোড় সংখ্যার দ্বিগুণ, বিশেষত $2 \cdot (2m+1)$। তৃতীয় সারির ক্রমযুগল $\tuple{2,m}$-গুলি এমন ক্রমস্থানে আছে, যার সংখ্যা একটি বিজোড় সংখ্যার চার গুণ, $4 \cdot (2m+1)$; এইভাবে চলতে থাকে। প্রত্যেক সারির জন্য $(2m+1)$-এর গুণকগুলি $1$, $2$, $4$, $8$, \dots{} হলো ঠিক~$2$-এর ঘাত: $1= 2^0$, $2 = 2^1$, $4 = 2^2$, $8 = 2^3$, \dots\@ বস্তুত, সংশ্লিষ্ট ঘাতাঙ্কটি সব সময় আলোচ্য ক্রমযুগলের প্রথম সদস্য। তাই ক্রমযুগল $\tuple{n,m}$-এর জন্য গুণক হলো $2^n$। এর থেকে সাধারণ সূত্র পাওয়া যায়: $2^n \cdot (2m+1)$। কিন্তু এটি ক্রমযুগল থেকে \emph{ধনাত্মক} পূর্ণসংখ্যায় একটি চিত্রণ, অর্থাৎ $\tuple{0,0}$-এর ক্রমস্থান~$1$। আমরা~$0$ ক্রমস্থান থেকে শুরু করতে চাইলে ফল থেকে~$1$ বিয়োগ করতে হবে। এতে পাই:
\end{explain}

\begin{ex}
নিচের সংজ্ঞা দ্বারা নির্ধারিত অপেক্ষক $h\colon \Nat^2 \to \Nat$:
\[
h(n,m) = 2^n (2m+1) - 1
\]
স্বাভাবিক সংখ্যার ক্রমযুগলের সেট~$\Nat^2$-এর একটি যুগলায়ন অপেক্ষক।
\end{ex}

\begin{explain}
সেই অনুযায়ী, আমাদের $\Nat^2$-এর দ্বিতীয় তালিকায়নে ক্রমযুগল $\tuple{0,0}$-এর সংকেত $h(0,0) = 2^0(2\cdot 0+1) - 1 = 0$;
$\tuple{1,2}$-এর সংকেত $2^{1} \cdot (2 \cdot 2 + 1) - 1 = 2
\cdot 5 - 1 = 9$; এবং $\tuple{2,6}$-এর সংকেত $2^{2} \cdot (2
\cdot 6 + 1) - 1 = 51$।
\end{explain}

কখনো কখনো স্বাভাবিক সংখ্যার ক্রমযুগল~$\Nat^2$-এর সংকেতায়ন করাই যথেষ্ট; সংকেতায়নটি সমাপতিত হওয়া আবশ্যক নয়। এই ধরনের সংকেতায়নের বিপরীত কেবল আংশিক অপেক্ষক হয়।

\begin{ex}
নিচের সংজ্ঞা দ্বারা নির্ধারিত অপেক্ষক $j\colon \Nat^2 \to \Nat^+$:
\[
j(n,m) = 2^n3^m
\]
হলো !!a{injective} অপেক্ষক $\Nat^2 \to \Nat$।
\end{ex}

\end{document}
